\section*{Overview}

Ternary search is the search pattern for one-dimensional optimization when the objective is \emph{unimodal}: it moves
in one direction until the optimum, then moves in the other direction after the optimum. The contest version is usually
``the function is convex / concave enough on one variable that I can compare two interior points and discard one side.''

\section*{When to Use It}

Use ternary search when:

\begin{itemize}
  \item the answer is the minimum or maximum of a function \(f(x)\),
  \item the domain is one-dimensional,
  \item \(f\) is unimodal on the searched interval,
  \item evaluating \(f(x)\) is much cheaper than enumerating every candidate.
\end{itemize}

Typical contest signals are: ``choose one coordinate'', ``best time / angle / height'', and ``the cost is convex after
the rest of the variables are fixed.''

\section*{Core Idea}

Pick two interior points \(m_1 < m_2\).

\begin{itemize}
  \item For minimization, if \(f(m_1) \le f(m_2)\), then the optimum cannot lie strictly to the right of \(m_2\).
  \item If \(f(m_1) > f(m_2)\), then the optimum cannot lie strictly to the left of \(m_1\).
\end{itemize}

So every comparison removes a constant fraction of the current interval.

\section*{Key Insight}

The search works because unimodality is stronger than ``there is one best point.'' It says the function has only one
direction change. Without that structural promise, comparing \(f(m_1)\) and \(f(m_2)\) does not tell you which side is
safe to discard.

\section*{Worked Problem}

\paragraph{Problem.}
There are \(n\) drones on a line. Drone \(i\) starts at position \(p_i\) and flies with speed \(s_i\). Choose a real
meeting point \(x\) that minimizes the time until \emph{all} drones can reach it.

\paragraph{Why ternary search fits.}
The objective is
\[
f(x) = \max_i \frac{|x - p_i|}{s_i}.
\]
Each term is a V-shaped convex function. The maximum of convex functions is still convex, so \(f(x)\) is unimodal for
minimization.

\paragraph{Algorithm outline.}

\begin{itemize}
  \item Search on \([\min p_i, \max p_i]\).
  \item Evaluate \(f(m_1)\) and \(f(m_2)\).
  \item Keep the half that still contains the minimum.
  \item Stop after enough iterations for the required precision.
\end{itemize}

\section*{Correctness Intuition}

For a convex function, once the value at \(m_1\) is already no larger than the value at \(m_2\), the right side has
started climbing or is not better than the left-middle region. So the minimizer cannot be hidden strictly beyond
\(m_2\). The symmetric argument handles the other case.

\section*{Complexity Analysis}

If one evaluation of \(f(x)\) costs \(T\), then:

\begin{itemize}
  \item continuous ternary search costs \(O(T \cdot I)\), where \(I\) is the chosen number of iterations,
  \item discrete ternary search costs \(O(T \log R)\) until the interval becomes small enough to brute-force.
\end{itemize}

\section*{Implementation}

The code includes:

\begin{itemize}
  \item a floating-point ternary search template,
  \item a sample evaluator for the drone meeting-time problem above.
\end{itemize}

For integer domains, I usually shrink with ternary search until only a handful of candidates remain, then test them all.

\section*{Common Pitfalls}

\begin{itemize}
  \item Using ternary search on a function that has multiple local optima.
  \item Applying it to integers without a final brute-force cleanup.
  \item Searching on a wider interval than the real feasible region and losing precision.
  \item Forgetting that many ``convex-looking'' formulas become non-unimodal after hidden constraints are introduced.
\end{itemize}

\section*{Variants / Extensions}

\begin{itemize}
  \item Discrete ternary search on arrays or DP transition indices.
  \item Binary search on the derivative sign when the derivative is monotone.
  \item Parametric reformulations where the optimization can be turned into a monotone decision problem instead.
\end{itemize}

\section*{Practice Problems}

\begin{itemize}
  \item Choose one real coordinate minimizing the worst travel time.
  \item Maximize a concave score depending on one continuous parameter.
  \item Optimize one split point after all other combinatorial choices are fixed.
\end{itemize}

\section*{References}

\begin{itemize}
  \item \href{https://cp-algorithms.com/num_methods/ternary_search.html}{cp-algorithms: Ternary Search}.
  \item \href{https://usaco.guide/gold/ternary-search}{USACO Guide: Ternary Search}.
  \item \href{https://codeforces.com/catalog}{Codeforces Catalog}.
\end{itemize}
